% !TeX program = xelatex
% Complete chapter-mode example for noteformyself 2.0.
\documentclass[sectionlevel=chapter,status=draft]{noteformyself}

% -----------------------------------------------------------------------------
% Chapter metadata and image cover
% -----------------------------------------------------------------------------
\title{A Chapter-Level Note with a Deliberately Long Explanatory Title}
\author{Tianle Yang}
\date{\today}
\coversentence{Geometry becomes clearer when the structure can breathe.}
\coverimage{placeholder.jpg}
\version{2.0.0}

% -----------------------------------------------------------------------------
% Tag context
% -----------------------------------------------------------------------------
\LoadNoteTagRegistry{tag-registry.example.tex}
\notesetup{
  book={Tag System Test Book},
  book-tag=0,
  chapter-tag=0CH,
  status=draft
}

\begin{document}

\maketitle
\tableofcontents

% =============================================================================
% A published section and subsection
% =============================================================================
\section[
  status=published,
  tag=0CH01,
  label=sec:chapter-preview
]{Published chapter material}

\begin{theorem}[
  title={A small test theorem},
  status=published,
  tag=0CH0101,
  label=thm:chapter-preview
]
  A quiet page can retain a strong hierarchy and a permanent identity.
\end{theorem}

\subsection[
  status=published,
  tag=0CH0102,
  label=subsec:chapter-preview
]{A tagged subsection}

The compact references are \cref{sec:chapter-preview},
\cref{subsec:chapter-preview}, and \cref{thm:chapter-preview}.

% =============================================================================
% Every theorem-style environment
% =============================================================================
\section[
  short-title={Environment gallery},
  status=draft,
  tag=0CH02,
  label=sec:chapter-gallery
]{Draft environment gallery}

\subsection[tag=0CH0201,label=subsec:chapter-gallery]{Statements and proofs}

\begin{slogan}
  A template should make the mathematical hierarchy visible at a glance.
\end{slogan}

\begin{mainthm}[title={Chapter comparison},label=mainthm:chapter]
  The three compilation modes share one statement vocabulary.
\end{mainthm}

\begin{definition}[title={Affine spectrum},label=def:chapter]
  For a ring \(A\), the space \(\operatorname{Spec}(A)\) consists of its prime
  ideals with the Zariski topology and its structure sheaf.
\end{definition}

\begin{proposition}[label=prop:chapter]
  The inverse image of a basic open set under a morphism of affine schemes is
  basic open.
\end{proposition}

\begin{lemma}[title={Localization test},label=lem:chapter]
  Localization commutes with finite colimits of modules.
\end{lemma}

\begin{corollary}[label=cor:chapter]
  A localization of a finitely presented module is finitely presented.
\end{corollary}

\begin{conjecture}[label=conj:chapter]
  The comparison remains true after replacing schemes by suitable stacks.
\end{conjecture}

\begin{question}[title={Minimal hypotheses},label=ques:chapter]
  Which finiteness assumption is actually used in the argument?
\end{question}

\begin{example}[tag=0CH0202,label=eg:chapter]
  The affine line is \(\mathbb{A}^1_k=\operatorname{Spec}(k[t])\). The displayed
  draft tag is provisional and is absent from \cref{eg:chapter}.
\end{example}

\begin{exercise}[label=ex:chapter]
  Compute the standard-open cover of \(\mathbb{P}^1_k\).
\end{exercise}

\begin{construction}[title={Gluing two affine lines},label=constr:chapter]
  Glue \(\operatorname{Spec}(k[t])\) and \(\operatorname{Spec}(k[t^{-1}])\)
  along \(\operatorname{Spec}(k[t,t^{-1}])\).
\end{construction}

\begin{notation}[label=not:chapter]
  Denote the resulting scheme by \(\mathbb{P}^1_k\).
\end{notation}

\begin{remark}[title={About draft references},label=rmk:chapter]
  Draft references display only the current type and number.
\end{remark}

\begin{proof}[Proof of \cref{mainthm:chapter}]
  Each mode installs the same public environments.

  \begin{step}[Compare counters]\label{step:chapter-counter}
    Ordinary statements share one counter; \texttt{mainthm} uses letters.
  \end{step}

  \begin{claim}[label=claim:chapter]
    The visual style does not alter the underlying amsthm semantics.
  \end{claim}

  \begin{case}[Book or chapter mode]\label{case:chapter-numbered}
    Statement numbers reset within sections.
  \end{case}

  \begin{case}[Section mode]\label{case:chapter-standalone}
    Statement numbers use a single local sequence.
  \end{case}
\end{proof}

The gallery supports references to \cref{def:chapter}, \cref{prop:chapter},
\cref{lem:chapter}, and \cref{cor:chapter}, together with
\cref{claim:chapter}, \cref{step:chapter-counter}, and
\cref{case:chapter-numbered}.

% =============================================================================
% Mathematics, diagrams, lists, and code
% =============================================================================
\subsection[tag=0CH0203,label=subsec:chapter-features]{Typesetting features}

The text font is Libertinus Serif, headings use Libertinus Sans, and source-like
material uses Source Code Pro. Mathematics uses Libertinus Math, with STIX Two
providing distinct calligraphic and script alphabets:
\[
  \mathcal{O}_X(1),\qquad
  \mathscr{H}^{i}(\mathscr{F}),\qquad
  \mathbb{P}^{n}_{k},\qquad
  \mathfrak{gl}(V),\qquad
  \mathbf{R}\!\operatorname{Hom}_{X}(E,F).
\]

\[
\begin{tikzcd}
  \operatorname{Spec}(B) \arrow[r] \arrow[d] & \operatorname{Spec}(A) \arrow[d] \\
  Y \arrow[r] & X
\end{tikzcd}
\]

\begin{enumerate}
  \item A first chapter-level task.
  \item A task with two parts.
    \begin{enumerate}
      \item establish the local claim;
      \item glue the local results.
    \end{enumerate}
\end{enumerate}

\begin{lstlisting}[
  language=Python,
  caption={A chapter-level code example},
  label=lst:chapter-example
]
def sections(chapter):
    return sorted(chapter.sections, key=lambda item: item.order)
\end{lstlisting}

See \cref{lst:chapter-example}. The draft structural reference
\cref{sec:chapter-gallery} omits its provisional tag.

\end{document}
